% !TeX program = xelatex
% ============================================================
% vkr/preamble.tex — преамбула ВКР ОП «Программная инженерия» ФКН
% НИУ ВШЭ. Воспроизводит оформление реальной ВКР 2025/2026
% (ГОСТ 7.32-2017): 12 pt, поля 30/15/20/20, интервал 1,5, абзац
% 1,25 см, сквозная нумерация рисунков/таблиц, «Оглавление».
%
% Данные о работе (ФИО, группа, год, руководитель…) приходят из
% common/meta.tex (\hse-макросы, перегенерируются перед сборкой) —
% не правьте их здесь, задавайте в настройках проекта.
% ============================================================

% ── Шрифты (автогенерируются из настроек проекта) ───────────
% Times New Roman 12 pt — как в реальной ВКР (ГОСТ 7.32). Семейства
% (основной/рубленый/моноширинный) задаются в настройках проекта и
% подставляются в common/fonts.tex ПЕРЕД каждой сборкой, каждый с
% независимым фолбэком на свободный клон с кириллицей (Tinos/DejaVu) —
% отсутствие шрифта не ломает сборку. Здесь НЕ правьте — меняйте в
% настройках проекта (устанавливайте шрифты в Settings → Шрифты).
% \hseDefaultMono — дефолт моношрифта ВКР, если он не выбран в настройках.
\newcommand{\hseDefaultMono}{Consolas}
\input{../common/fonts}
% Язык оформления по project.lang. §2.5.3.1: англоязычная ВКР пишется на
% английском (главный язык — english); реферат остаётся двуязычным (Реферат +
% Abstract) в обоих случаях. Русский — язык по умолчанию.

\usepackage[russian,english]{babel}
\addto\captionsenglish{\renewcommand{\bibname}{References}}

\usepackage{url}
\usepackage[a4paper,left=30mm,right=15mm,top=20mm,bottom=20mm]{geometry}
\usepackage{setspace}
\onehalfspacing
\usepackage{indentfirst}
\setlength{\parindent}{1.25cm}
\setlength{\parskip}{6pt}

\usepackage{microtype}
\usepackage{fvextra}
\usepackage{csquotes}
\usepackage{graphicx}
\usepackage{letltxmacro}
\LetLtxMacro{\oldincludegraphics}{\includegraphics}
\renewcommand{\includegraphics}[2][]{%
  \IfFileExists{#2}{%
    \oldincludegraphics[#1]{#2}%
  }{%
    \IfFileExists{#2.png}{%
      \oldincludegraphics[#1]{#2.png}%
    }{%
      \IfFileExists{#2.jpg}{%
        \oldincludegraphics[#1]{#2.jpg}%
      }{%
        \fcolorbox{red}{white}{%
          \begin{minipage}{0.5\textwidth}
            \centering
            \textcolor{red}{\textbf{Файл не найден:}} \\
            \texttt{\detokenize{#2}}
          \end{minipage}%
        }%
      }%
    }%
  }%
}
\usepackage{xcolor}
\usepackage{hyperref}
\usepackage{lastpage}
\usepackage{float}
\usepackage{array}
\usepackage{booktabs}
\usepackage{longtable}
\usepackage{adjustbox}
\usepackage{xurl} % Allow long URLs to break
\urlstyle{same}
\usepackage{hyphenat} % Better hyphenation
\usepackage{caption}
\usepackage{cite} % Better citation handling
\DeclareCaptionLabelSeparator{emdash}{~---~}
\captionsetup[table]{position=top,labelsep=emdash,justification=raggedright,singlelinecheck=false}
\captionsetup[figure]{position=bottom,labelsep=emdash}

% --- Настройка подсветки кода ---
\usepackage{listings}

\renewcommand{\lstlistingname}{Listing}
\renewcommand{\lstlistlistingname}{List of Listings}

\definecolor{codegreen}{rgb}{0,0.6,0}
\definecolor{codegray}{rgb}{0.5,0.5,0.5}
\definecolor{codepurple}{rgb}{0.58,0,0.82}
\definecolor{backcolour}{rgb}{0.95,0.95,0.92}
\definecolor{keycolor}{rgb}{0.13,0.29,0.53}
\definecolor{stringcolor}{rgb}{0.31,0.60,0.02}

((* raw *))
% raw-блок: в определении ниже есть литеральные двойные фигурные скобки
% (listings literate), которые рендер иначе принял бы за начало Jinja-
% выражения -> TemplateSyntaxError, и весь файл скопировался бы как есть
% (языковые if-ветки не сработали бы). raw защищает эти скобки.
\lstdefinelanguage{yaml}{
  keywords={true,false,null,y,n},
  keywordstyle=\color{darkgray}\bfseries,
  basicstyle=\small\ttfamily,
  breaklines=true,
  comment=[l]{\#},
  commentstyle=\color{codegray}\itshape,
  stringstyle=\color{stringcolor},
  moredelim=[l][\color{keycolor}\bfseries]{:},
  morestring=[b]',
  morestring=[b]",
  literate={-}{{\textcolor{keycolor}{-}}}1
}
((* endraw *))

\lstnewenvironment{CodeBlock}[1][]
{
    \lstset{
        language=yaml,
        backgroundcolor=\color{backcolour},
        commentstyle=\color{codegreen},
        keywordstyle=\color{magenta},
        numberstyle=\tiny\color{codegray},
        stringstyle=\color{codepurple},
        basicstyle=\footnotesize\ttfamily,
        breakatwhitespace=false,
        breaklines=true,
        captionpos=b,
        keepspaces=true,
        numbers=left,
        numbersep=5pt,
        showspaces=false,
        showstringspaces=false,
        showtabs=false,
        tabsize=2,
        frame=single,
        rulecolor=\color{gray!30},
        xleftmargin=10pt,
        xrightmargin=10pt,
        #1
    }
}{}
% -------------------------------

\usepackage{enumitem}
\setlist{nosep}
\setlist[itemize,1]{label=--}
\usepackage{chngcntr}
\counterwithout{table}{chapter}
\counterwithout{figure}{chapter}
\renewcommand{\thetable}{\arabic{table}}
\renewcommand{\thefigure}{\arabic{figure}}
\usepackage{fancyhdr}
\pagestyle{plain}

\usepackage{titlesec}
\usepackage{tikz}
\usetikzlibrary{shapes.geometric, arrows, positioning, calc, er, fit}
\tikzstyle{startstop} = [rectangle, rounded corners, minimum width=3cm, minimum height=1cm,text centered, draw=black, fill=red!30]
\tikzstyle{io} = [trapezium, trapezium left angle=70, trapezium right angle=110, minimum width=3cm, minimum height=1cm, text centered, draw=black, fill=blue!30]
\tikzstyle{process} = [rectangle, minimum width=3cm, minimum height=1cm, text centered, draw=black, fill=orange!30]
\tikzstyle{decision} = [diamond, minimum width=3cm, minimum height=1cm, text centered, draw=black, fill=green!30]
\tikzstyle{arrow} = [thick,->,>=stealth]
\tikzstyle{box} = [rectangle, draw, text centered, minimum height=1cm, minimum width=2.5cm, rounded corners]
\tikzstyle{database} = [cylinder, draw, shape border rotate=90, aspect=0.25, text centered, minimum height=1.5cm, minimum width=1.2cm]
\tikzstyle{entity} = [rectangle, draw, text centered, minimum height=1cm, minimum width=2.5cm, fill=blue!10]
\tikzstyle{relationship} = [diamond, draw, text centered, minimum height=1cm, minimum width=1.5cm, fill=red!10]

\emergencystretch=5em % Help with overfull hboxes
\tolerance=2000
\hyphenpenalty=300
\exhyphenpenalty=300
\titleformat{\chapter}{\bfseries\centering}{\thechapter}{1em}{}
\titleformat{\section}{\bfseries}{\thesection}{1em}{}
\titlespacing*{\section}{1.25cm}{*2}{6pt}
\titleformat{\subsection}{\bfseries}{\thesubsection}{1em}{}
\titlespacing*{\subsection}{1.25cm}{*2}{6pt}
\titleformat{\subsubsection}{\bfseries}{\thesubsubsection}{1em}{}


\renewcommand{\contentsname}{Contents}

\setcounter{tocdepth}{3}
\setcounter{secnumdepth}{3}

\hypersetup{unicode=true,breaklinks=true,bookmarksopen=true,bookmarksnumbered=true,colorlinks=true,linkcolor=black,urlcolor=blue,citecolor=blue}

\newcommand{\cmpplus}{\textbf{+}}
\newcommand{\cmpminus}{\textbf{--}}
\newcommand{\cmpstar}{\textbf{+*}}
\newcommand{\cmpplanned}{\textcolor{blue}{$\pm$}}

% Сокращения для названий аналогов в матрицах сравнения с гиперссылками на официальные сайты
\newcommand{\anYandexCloud}{\href{https://yandex.cloud/}{Yandex Cloud}}
\newcommand{\anAWS}{\href{https://aws.amazon.com/}{AWS}}
\newcommand{\anGoogleCloud}{\href{https://cloud.google.com/}{Google Cloud}}
\newcommand{\anAzure}{\href{https://azure.microsoft.com/}{Azure}}
\newcommand{\anOpenShift}{\href{https://www.redhat.com/en/technologies/cloud-computing/openshift}{OpenShift}}
\newcommand{\anCloudRu}{\href{https://cloud.ru/}{Cloud.ru}}
\newcommand{\anPortainer}{\href{https://www.portainer.io/}{Portainer}}
\newcommand{\anMogenius}{\href{https://mogenius.com/}{Mogenius}}
\newcommand{\anHeroku}{\href{https://www.heroku.com/}{Heroku}}
\newcommand{\anPlatformSh}{\href{https://platform.sh/}{Platform.sh}}
\newcommand{\anHumanitec}{\href{https://humanitec.com/}{Humanitec}}
\newcommand{\anMiaPlatform}{\href{https://mia-platform.eu/}{Mia Platform}}
\newcommand{\anExoscale}{\href{https://www.exoscale.com/}{Exoscale}}
\newcommand{\anQovery}{\href{https://www.qovery.com/}{Qovery}}
\newcommand{\anAppvia}{\href{https://www.appvia.io/}{Appvia}}
\newcommand{\anNullstone}{\href{https://nullstone.io/}{Nullstone}}
\newcommand{\anNagios}{\href{https://www.nagios.org/}{Nagios}}
\newcommand{\anPorter}{\href{https://porter.run/}{Porter}}
\newcommand{\anCoherence}{\href{https://withcoherences.webflow.io/}{Coherence}}
\newcommand{\anDeployMe}{\href{https://deployme.app/}{Deploy Me}}

% Трассировка требований (панель «Трассировка требований»):
%   \req{ID}{Описание} — определение (в ТЗ); \reqref{ID} — ссылка (в ВКР и др.).
\newcommand{\req}[2]{\phantomsection\label{req:#1}\textbf{#1.}\enspace #2}
\newcommand{\reqref}[1]{\mbox{(треб.~#1)}}

% Форматирование номеров в списке литературы с точкой
\makeatletter
\renewcommand\@biblabel[1]{#1.}
\makeatother

% ── Данные проекта (автогенерируемые макросы \hse…) ─────────
\input{../common/meta}
